\begin{center}
\hypertarget{native-052-article035}{}\pdfbookmark[2]{35.}{article-035}35.
\end{center}
Expressiones analyticas trium coëfficientium primorum seriei secundae $w^{\prime}$, $w^{\prime\prime}$, $w^{\prime\prime\prime}$ \textsc{Bessel} noster evolvit, simulque pro valoribus numericis duorum primorum $w^{\prime}$, $w^{\prime\prime}$ tabulam ad singulos argumenti $w$ gradus constructam addidit (\textsc{von Zach} Monatliche Correspondenz, vol. XII. p. 197). Pro coëfficiente primo $w^{\prime}$ tabula iam ante habebatur a \textsc{Simpson} computata, quae operi clar. \textsc{Olbers} supra laudato annexa est. In plerisque casibus hacce methodo adiumento tabulae \textsc{Bessel}ianae anomaliam veram e tempore praecisione sufficiente \SourceLeaf{053}determinare licet: quod adhuc desiderandum relinquitur, ad haecce fere momenta reducitur:

I. In problemate inverso, temporis puta ex anomalia vera determinatione, ad methodum quasi indirectam confugere atque $w$ ex $v$ tentando derivare oportet. Cui incommodo ut obveniretur, series prior eodem modo tractata esse deberet ac secunda: et quum facile perspiciatur, $-v^{\prime}$ esse eandem functionem ipsius $v$, qualis $w^{\prime}$ est ipsius $w$, ita ut tabula pro $w^{\prime}$ signo tantum mutato pro $v^{\prime}$ inservire possit, nihil tam requireretur nisi tabula pro $v^{\prime\prime}$, quo utrumque problema aequali praecisione solvere liceat.

II. Interdum utique occurrere possunt casus, ubi excentricitas ab unitate parum quidem differt, ita ut methodi generales supra expositae praecisionem haud sufficientem dare videantur, nimis tamen etiamnum, quam ut in methodo peculiari modo adumbrata effectum potestatis tertiae ipsius $\delta$ altiorumque tuto contemnere liceat. In motu imprimis hyperbolico eiusmodi casus sunt possibiles, ubi, sive illas methodos adoptes sive hanc, errorem plurium secundorum evitare non possis, siquidem tabulis vulgaribus tantum ad septem figuras constructis utaris. Etiamsi vero huiusmodi casus in praxi raro occurrant, aliquid certe deesse videri posset, si in \emph{omnibus} casibus anomaliam veram intra $0^{\prime\prime}{,}1$ aut saltem $0^{\prime\prime}{,}2$ determinare non liceret, nisi tabulae maiores consulerentur, quas tamen ad libros rariores referendas esse constat. Haud igitur prorsus superfluam visum iri speramus expositionem methodi peculiaris, qua iamdudum usi sumus, quaeque eo etiam nomine se commendabit, quod ad excentricitates ab unitate parum diversas haud limitata est, sed hocce saltem respectu applicationem generalem patitur.

\begin{center}
\hypertarget{native-053-article036}{}\pdfbookmark[2]{36.}{article-036}36.
\end{center}
Antequam hanc methodum exponere aggrediamur, observare conveniet, incertitudinem methodorum generalium supra traditarum in orbitis ad parabolae similitudinem vergentibus sponte desinere, simulac $E$ vel $F$ ad magnitudinem considerabilem increverint, quod quidem in magnis demum a Sole distantiis fiet. Quod ut ostendámus, errorem maximum in ellipsi possibilem, quem in art. 32, IV invenimus $\frac{3\omega ea\sin v}{\lambda r}\cdot206265^{\prime\prime}$ ita exhibemus
\[
\frac{3\omega e\BareRoot{(1-ee)}.\sin E}{\lambda(1-e\cos E)^2}\cdot206265^{\prime\prime},
\]
\SourceLeaf{054}unde sponte patet, errorem arctis semper limitibus circumscriptum esse, simulac $E$ valorem considerabilem acquisiverit, sive simulac $\cos E$ ab unitate magis recesserit quantumvis magna sit excentricitas. Quod adhuc luculentius apparebit per tabulam sequentem, in qua valorem numericum maximum istius formulae pro quibusdam valoribus determinatis computavimus (pro septem figuris decimalibus):
\begin{center}
\begin{tabular}{r|lr}
$E=10^0$ & error maximus $=$ & $3^{\prime\prime}{,}04$\\
$20$ & & $0{,}76$\\
$30$ & & $0{,}34$\\
$40$ & & $0{,}19$\\
$50$ & & $0{,}12$\\
$60$ & & $0{,}08$
\end{tabular}
\end{center}
\noindent Simili modo res se habet in hyperbola, ut statim apparet, si expressio in art. 32, VII eruta sub hanc formam ponitur $\frac{\omega\cos F.(\cos F+3e\sin F)\BareRoot{(ee-1)}}{\lambda(e-\cos F)^2}206265^{\prime\prime}$. Valores maximos huius expressionis pro quibusdam valoribus determinatis ipsius $F$ tabula sequens exhibet:
\begin{center}
\begin{tabular}{r|r|r|r}
\hline
$F$ & \multicolumn{2}{c|}{$u$} & error maximus\\
\hline
$10^0$ & $1{,}192$ & $0{,}839$ & $8^{\prime\prime}{,}66$\\
$20$ & $1{,}428$ & $0{,}700$ & $1{,}38$\\
$30$ & $1{,}732$ & $0{,}577$ & $0{,}47$\\
$40$ & $2{,}144$ & $0{,}466$ & $0{,}22$\\
$50$ & $2{,}747$ & $0{,}364$ & $0{,}11$\\
$60$ & $3{,}732$ & $0{,}268$ & $0{,}06$\\
$70$ & $5{,}671$ & $0{,}176$ & $0{,}03$
\end{tabular}
\end{center}
\noindent Quoties itaque $E$ vel $F$ ultra $40^0$ vel $50^0$ egreditur (qui tamen casus in orbitis a parabola parum discrepantibus haud facile occurret, quum corpora coelestia in talibus orbitis incedentia in tantis a Sole distantiis oculis nostris plerumque se subducant), nulla aderit ratio, cur methodum generalem deseramus. Ceterum in tali casu etiam series de quibus in art. 34 egimus nimis lente convergerent: neutiquam igitur pro defectu methodi nunc explicandae haberi potest, quod iis imprimis casibus adaptata est, ubi $E$ vel $F$ ultra valores modicos nondum excrevit.
